% ==================================================================
% id:       lim-004
% topic:    גבולות ורציפות של פונקציות
% subtopic: חישוב גבולות של פונקציות
% solution: full
% ==================================================================
\textbf{
בכל אחד מהסעיפים הבאים,
חשבו את הגבול:}

\begin{enumerate}\bMs{2}
\item
$\ds{
\lim_{x\to -3}\frac{x^2+6x+9}{3x^2+7x-6}
} $
\ifsol
\begin{solbox}
לאחר ההצבה של
$x=-3$
בפונקציה נקבל ביטוי
$"\tfrac{0}{0}"$,
כלומר, לא נוכל להיעזר באריתמטיקה ולהסיק מיד מה גבולה של הפונקציה כש-
$x$
שואף ל-
$(-3)$.
בגלל שבמונה ובמכנה יש פולינומים המתאפסים בנקודה
$-3$
אז נצליח להציגם בצורה הבאה
$$,f(x)=\frac{x^2+6x+9}{3x^2+7x-6}=\frac{(x-(-3))^2}{(x-(-3))(3x-2)}=\frac{x+3}{3x-2}$$
לכן
\begin{align*}
\ds{
\lim_{x\to -3}\frac{x^2+6x+9}{3x^2+7x-6}
}
= \ds{
\lim_{x\to -3}\frac{x+3}{3x-2}}
\overset{\ast}{=}\ \frac{-3+3}{3\cdot(-3)-2}
= 0.
\end{align*}
$(\ast)$
יש כאן שימוש באריתמטיקה של גבולות של פונקציות.
\end{solbox}
\fi

\item
$\ds{
\lim_{x\to 4} \left( \frac{1}{4-x} \ - \  \frac{48}{64-x^3} \right)
} $
\ifsol
\hint{פרקו את $64-x^3$ בנוסחת הפרש חזקות שלישיות ועברו למכנה משותף.}
\begin{solbox}
הצבה של
$x=4$
בפונקציה נותנת ביטוי מהצורה
$"\infty-\infty"$,
ביטוי שאנו לא יכולים להסיק מסקנה מידית ממנו.
נשים לב כי
$$,64-x^3=4^3-x^3=(4-x)(4^2+4x+x^2)$$
לכן
נציג
\begin{align*}
    \frac{1}{4-x}-\frac{48}{64-x^3}&=\ \frac{1}{4-x}-\frac{48}{(4-x)(16+4x+x^2)}\\
    &=\ \frac{16+4x+x^2-48}{(4-x)(16+4x+x^2)}\\
    &=\ \frac{x^2+4x-32}{(4-x)(16+4x+x^2)}\\
    &=\ \frac{(x-4)(x+8)}{(4-x)(16+4x+x^2)}\\
    &=\ -\frac{x+8}{16+4x+x^2},
\end{align*}
מכאן
נשתמש באריתמטיקה של גבולות ונקבל
\begin{align*}
   \lim_{x\to 4}\left(\frac{1}{4-x}-\frac{48}{64-x^3}\right)&=\ -\lim_{x\to 4}\frac{x+8}{16+4x+x^2}\\
   &=\ -\frac{4+8}{16+4\cdot4+4^2}
   = -\frac{1}{4}.
\end{align*}
\end{solbox}
\fi

\item
$\ds{
\lim_{x\to 5} \frac{ \sqrt{x-1} - 2}{x^2-25}
} $
\ifsol
\hint{במונה יש הפרש עם שורש: הכפילו וחלקו בצמוד.}
\begin{solbox}
הצבה של
$x=5$
בפונקציה נותנת לנו
$"\frac{0}{0}"$.
הביטוי במכנה ניתן לכתיבה

$x^2-25=(x-5)(x+5)$.
הדחף הראשון להפטר מהפרש של שני שורשים הוא כפל בצמוד, זה עשוי לעזור כאן
\begin{align*}
\frac{\sqrt{x-1}-2}{x^2-25}&=\ \frac{\sqrt{x-1}-2}{(x-5)(x+5)}\cdot{\color{red}\frac{\sqrt{x-1}+2}{\sqrt{x-1}+2}}\\
&=\ \frac{x-5}{(x-5)(x+5)(\sqrt{x-1}+2)}\\
&=\ \frac{1}{(x+5)(\sqrt{x-1}+2)},
\end{align*}
מכאן
\begin{align*}
    \lim_{x\to 5} \frac{ \sqrt{x-1} - 2}{x^2-25}&=\ \lim_{x\to 5} \frac{1}{(x+5)(\sqrt{x-1}+2)}\\
    &=\ \frac{1}{(5+5)(\sqrt{5-1}+2)}\\
    &=\ \frac{1}{40}
\end{align*}
הערה: בחישוב הגבול השתמשנו באריתמטיקה של גבולות: הגבול של סכום, של מכפלה ושל מנה של פונקציות מתכנסות שווה לסכום, למכפלה ולמנה של הגבולות, כל עוד הגבול של המכנה שונה מ-$0$.
\end{solbox}
\fi

\item $\ds{ \lim_{x \to 4} \lr{\frac{x + 3}{x - 4}-\frac{49}{x^2 - x - 12}}}$
\ifsol
\begin{solbox}
הצבה של
$x=4$
בפונקציה נותנת ביטוי
מהצורה
$"\infty-\infty"$.
נציג
\begin{align*}
\frac{x + 3}{x - 4}-\frac{49}{x^2 - x - 12}&=\ \frac{x + 3}{x - 4}-\frac{49}{(x-4)(x+3)}\\
&=\ \frac{x^2+6x-40}{(x-4)(x+3)}\\
&=\ \frac{(x-4)(x+10)}{(x-4)(x+3)}
=\frac{x+10}{x+3},
\end{align*}
מכאן
\begin{align*}
    \lim_{x \to 4} \lr{\frac{x + 3}{x - 4}-\frac{49}{x^2 - x - 12}}&=\ \lim_{x \to 4}\frac{x+10}{x+3}
    = \frac{4+10}{4+3}
    = 2.
\end{align*}
הערה: בחישוב הגבול השתמשנו באריתמטיקה של גבולות: הגבול של סכום, של מכפלה ושל מנה של פונקציות מתכנסות שווה לסכום, למכפלה ולמנה של הגבולות, כל עוד הגבול של המכנה שונה מ-$0$.
\end{solbox}
\fi
\eMs
\end{enumerate}
